\documentclass[10pt,a4paper]{article}
\usepackage[margin=1in]{geometry}

\usepackage{afterpage}
\usepackage{diagbox}
\usepackage[UTF8,heading = true]{ctex}


%%%% 宏包 %%%%%%%%%%%
\usepackage{amsmath,amssymb,amsfonts,mathrsfs}
\usepackage{graphicx}
\usepackage{ulem}
\usepackage{pdfpages}
\usepackage{lipsum}
\usepackage{booktabs}

\usepackage{array,tabularx,longtable}
	\newcolumntype{Y}{>{\centering\arraybackslash}X}
\usepackage{multirow}
\usepackage{listings}
\usepackage{enumerate }

\usepackage{tikz}
\usepackage{tikz-qtree,tikz-qtree-compat}
	\usetikzlibrary{arrows}

\usepackage{siunitx}
\usepackage{anyfontsize}

%\usepackage{fontspec}
\usepackage{xeCJK}

%\setCJKmainfont{SimSun.TTC}

\usepackage{pdfpages}


\usepackage{setspace}
\usepackage{color}
\usepackage{titletoc}

%解决参考文献的问题
\usepackage[comma]{natbib}
\setcitestyle{numbers,square}

\usepackage[explicit]{titlesec}
\titlespacing{\section}{0pt}{0pt plus .0pt }{20pt plus .0pt}
%\titlespacing{\section*}{0pt}{15pt plus .0pt}{19pt plus .0pt}
\titlespacing{\subsection}{0pt}{6pt plus .0pt}{6pt plus .0pt}


%\renewcommand{heiti}{SimHei.TTF}          % 黑体



\usepackage{algorithm,algorithmic}%写算法的环境

\usepackage{float}%解决表格与文字分开显示的问题



\ctexset{
    %contentsname = {\zihao{3}\heiti\bfseries{目录}},
%	figurename = {Figure},
%	tablename = {Table},
%	bibname = {References},
    section={
        name={第,章},
        format=\bfseries\centering\heiti\zihao{3},
        numberformat=\bfseries\heiti\zihao{3},
            },
    subsection={
        format=\bfseries\heiti\zihao{4},
        numberformat=\bfseries\heiti\zihao{4}
    },
    subsubsection={
        format=\bfseries\heiti\zihao{-4}
    }
        }



















%%%% 页面设置 %%%%%%%%%%

\usepackage[framed,numbered,autolinebreaks,useliterate]{mcode}%用来写代码的库
\usepackage{fancyhdr}
\pagestyle{fancy}


%--目录页眉页脚

\fancypagestyle{yemei}
{
    \fancyhf{}
    \fancyhead[C]{\zihao{5}\songti{\nouppercase{\leftmark}}}
    \fancyfoot[c]{\thepage}
    \renewcommand{\headrulewidth}{0.5pt}
    %\renewcommand{\footrulewidth}{1.5pt} %给页脚加上一条横线，在论文中没有使用的价值
}
\fancypagestyle{mulu}
{
    \fancyhf{}
    \fancyhead[C]{\zihao{5}\songti{目录}}
    \renewcommand{\headrulewidth}{0.5pt}
}

%改变颜色

\usepackage[colorlinks,linkcolor=blue,citecolor=blue,bookmarks=true,bookmarksnumbered=true]{hyperref}


% 目录字体设置
\usepackage{subfigure}
\usepackage[subfigure]{tocloft}      %必须这么写，否则会报错

\renewcommand{\baselinestretch}{1.2} %调整行距为1.2倍标准行距
\renewcommand{\cftsecleader}{\cftdotfill{0.6}}
\renewcommand{\cftsubsecleader}{\cftdotfill{0.6}}
\renewcommand{\cftsubsubsecleader}{\cftdotfill{0.6}}
\renewcommand{\cftsecfont}{\zihao{-4}\songti}    %设置section条目的字体
\renewcommand{\cftsubsecfont}{\zihao{-4}\songti} %设置subsection条目的字体
\renewcommand{\cftsubsubsecfont}{\zihao{-4}\songti} %设置subsection条目的字体

%解决参考文献是蓝色的问题，让参考不显示颜色。
\hypersetup{
    %colorlinks=true,
    %linkcolor=blcak,
    filecolor=black,      
    urlcolor=black,
}


\setlength\baselineskip{20pt}%调整正文内容的行距
\setlength{\headheight}{15pt}
\let\heiti\relax % 清除旧定义，
\let\songti\relax

\usepackage{caption}


\usepackage{threeparttable}    %给表格添加注释
%%%%%%%%%%%%%%%%%%%%%%%%%%%%%%%%%%%%%%%%%%%%%%%%%%%%%%%%%%%%
\begin{document}
    \bibliographystyle{plainnat}
    \includepdf{cover.pdf} %调用封面.
	
    \renewcommand\thetable{\thesection.\arabic{table}} 
    \renewcommand\thefigure{\thesection.\arabic{figure}} %解决图表的编号问题
    
    \newCJKfontfamily{\heiti}[AutoFakeBold={1.7}]{SimHei.TTF}
    %\newCJKfontfamily{\songti}[AutoFakeBold={3}]{SimSun.TTC}

	
    \renewcommand{\figurename}{图}
    \renewcommand\tablename{表}%将table改成表
    \captionsetup{labelformat=default,labelsep=default}%将表格加的那个小点取消
    \renewcommand{\thesubfigure}{\empty}%将subfigure下面的a,b,c给取消掉


    %% 中文摘要
    %\clearpage
    \fancyhead[C]{\zihao{5}\songti{摘要}}
    \phantomsection
    \addcontentsline{toc}{section}{摘要}

    
    \vspace*{0pt}
    \centerline{\zihao{3}\heiti\bfseries 摘\quad 要}
    \vspace*{16pt}

    \linespread{1.4}
    \zihao{-4} 

    当蜘蛛网无情地查封了我的炉台
　　  当灰烬的余烟叹息着贫困的悲哀
　　  我依然固执地铺平失望的灰烬
　　  用美丽的雪花写下：相信未来
　　  当我的紫葡萄化为深秋的露水
　　  当我的鲜花依偎在别人的情怀
　　  我依然固执地用凝霜的枯藤
　　  在凄凉的大地上写下：相信未来
　　  我要用手指那涌向天边的排浪
　　  我要用手撑那托住太阳的大海
　　  摇曳着曙光那枝温暖漂亮的笔杆
　　  用孩子的笔体写下：相信未来
　　  我之所以坚定地相信未来
　　  是我相信未来人们的眼睛
　　  她有拨开历史风尘的睫毛
　　  她有看透岁月篇章的瞳孔
　　  不管人们对于我们腐烂的皮肉
　　  那些迷途的惆怅、失败的苦痛
　　  是寄予感动的热泪、深切的同情
　　  还是给以轻蔑的微笑、辛辣的嘲讽
　　  我坚信人们对于我们的脊骨
　　  那无数次的探索、迷途、失败和成功
　　  一定会给予热情、客观、公正的评定
　　  是的，我焦急地等待着他们的评定
　　  朋友，坚定地相信未来吧
　　  相信不屈不挠的努力
　　  相信战胜死亡的年轻
　　  相信未来、热爱生命



    \bigskip

    \noindent{\zihao{-4}\songti \textbf{关键词:}}
    相信未来，热爱生命
    \pagenumbering{Roman}%使用罗马字体



    \newpage
    \phantomsection\addcontentsline{toc}{section}{Abstract}%将一些非章节加到目录之中
    \fancyhead[C]{\zihao{5}\songti{Abstract}}
    \vspace*{0pt}
    \centerline{\zihao{3}\bfseries Abstract}
    \vspace*{20pt}

    \linespread{1.4}\zihao{-4}
    ***************中文摘要的翻译即可****************
   

    \bigskip
    \noindent\textbf{\zihao{-4} Keywords:} 
    believe in yourself always



    \newpage
    \setcounter{secnumdepth}{4}%这个是让你的目录在正文中也不显示
    \setcounter{tocdepth}{2}%设定目录深度为2，即只显示到二级标题为止

    \renewcommand\contentsname{\zihao{3}\heiti\textbf{目\quad 录}}
  
    \begin{spacing}{1.7}
        {\centering\tableofcontents} %列出目录  
    \end{spacing}
    \thispagestyle{mulu}
    \vspace*{16pt}
	




    \setcounter{page}{0} %将本页记为第0页，重新编号
    %\thispagestyle{empty}%取消本页的页眉页脚

    \pagestyle{yemei}
    \newpage
    \pagenumbering{arabic}%重新设置成为阿拉伯数字
    \setcounter{page}{1} %将本页记为第1页，重新编号

    \section{绪论}
    %此命令放在正文，紧跟\section{}命令，手动为页眉添加节标题
    \markboth{\text{第}\thesection\text{章} \quad 绪论}{\text{第}\thesection\text{章} \quad 绪论}   

    \subsection{研究背景与研究目的}
    
        北国风光，千里冰封，万里雪飘。
      望长城内外，惟余莽莽;大河上下，顿失滔滔。
     山舞银蛇，原驰蜡象，欲与天公试比高。
     须晴日，看红装素裹，分外妖娆。
      江山如此多娇，引无数英雄竞折腰。
      惜秦皇汉武，略输文采;唐宗宋祖，稍逊风骚。
     一代天骄，成吉思汗，只识弯弓射大雕。
     俱往矣，数风流人物，还看今朝。
    


    

    \subsubsection{放自己的标题}
    $\times \times \times\times \times \times\times \times \times\times \times \times\times \times \times\times \times \times\times \times \times$

    \subsubsection{$\times \times \times\times \times \times$}
    $\times \times \times\times \times \times\times \times \times\times \times \times\times \times \times\times \times \times\times \times \times$

    \subsection{大括号里面放标题}

    





    \newpage
    \section{公式的编辑方法}
    %此命令放在正文，紧跟\section{}命令，手动为页眉添加节标题
    \markboth{\text{第}\thesection\text{章} \quad 公式的编辑方法}{\text{第}\thesection\text{章} \quad 公式的编辑方法}  


    \subsection{公式的输入方法}
    %%%%第一种：$$与$$$$
    使用美元符号进行行内公式输入：$a^2+b^2=c^2$.%代表是输入行内公式

    使用双美元符号进行行间公式输入：$$a^2+b^2=c^2.$$%代表是输入行间公式

    使用math环境进行行间公式输入：%采用math环境，是行内公式
    \begin{math}
        a^2+b^2=c^2.
    \end{math}

    %%%%第二种：equation环境
    equation环境下启用自动编号：
    \begin{equation}%不加星号表示启用自动编号
        x^3+x^2+x+1=0
    \end{equation}

    equation*环境下不启用自动编号，手动编号：
    \begin{equation*}%加星号表示不启用自动编号，使用\tag{}可以实现手动编号
        x^3+x^2+x+1=0
        \tag{2.1}
    \end{equation*}

    \subsection{复杂公式的输入方法}
    本文只介绍使用equation环境的输入方法，这种方法可以应对大部分复杂公式。例：
    \begin{equation}
    \label{f1}%给公式进行赋名，方便后面引用
        \left\{ %这种方法可以打出左括号
        \begin{aligned}
            &\phi _t =-(\phi^3 -g\phi ^2-\epsilon \phi +(1+\Delta  )^2 \phi),\\
            &\phi(x,y,0)=\phi_0(x,y),\\
            &\phi (x+L,y,t)=\phi(x,y,t),\\
            &\phi (x,y+L,t)=\phi (x,y,t),
        \end{aligned}
		\right.\tag{2.2}
    \end{equation}

    如果想要引用上面的公式，可以使用ref命令。例如：公式$(\ref{f1})$是一个复杂的数学公式。

    \subsection{参考文献的引用}
    如果想要在正文里面引用参考文献，刚可以使用cite命令。例如：参考文献$\cite{ref1}$是一个很重要的文献。在pdf文件中点击可以跳转到对应的参考文献。









    \newpage
    \section{图表的使用}
    在本节中，将具体对图和表的引用进行说明。

    %此命令放在正文，紧跟\section{}命令，手动为页眉添加节标题
    \markboth{\text{第}\thesection\text{章} \quad 图表的使用}{\text{第}\thesection\text{章} \quad 图表的使用}  


    \subsection{图的引用}
    LaTex提供自己的画图工具包，应该是tikz.这个画出来就是矢量图，但学起来相对复杂一点，故不做介绍。本节只对如何使用标量图进行介绍。

    第一步：把你需要的图片放到和main.tex一个位置。

    第二步：调用图片。第一张图如何调用如下：

    \begin{figure}[htbp]
        \centering
        \includegraphics[height=6cm,width=7.5cm]{0.jpg} 
        \caption{\songti{这是一张图的调用.}}
        \captionsetup{font={small,bf,stretch=1}}
        \label{figure1}
    \end{figure}

    多张图如何调用如下：
    \begin{figure}[htbp]
        \centering
        \subfigure[$T=0$]{
        \includegraphics[height=4cm,width=4.5cm]{0.jpg} 
        }
        \quad
        \subfigure[$T=10$]{
        \includegraphics[height=4cm,width=4.5cm]{10.jpg} 
        }
        \quad
        \subfigure[$T=20$]{
        \includegraphics[height=4cm,width=4.5cm]{20.jpg} 
        }
        \caption{\songti{这是多张图的调用}}
        \captionsetup{font={small,bf,stretch=1}}
        \label{figure2}
    \end{figure}

    依然是采用ref命令来引用图。例如：图$\ref{figure1}$是一张图，而图$\ref{figure2}$是多张图。
    

    \subsection{表的生成}
    表格在latex中是可以生成的。我们讲3个例子。

    例1：
    \begin{table}[htbp]
        \centering
        \setlength{\tabcolsep}{10mm}
        \begin{tabular}{ccc}  
            \hline
            年份&工资(元)&涨幅\\ \hline
            $1900$&$10$&--\\ 
            $2000$&$100$&$1000\%$\\ 
            $2100$&$1000$&$1000\%$\\  \hline
        \end{tabular} 
        \caption{\centering{简单的表格}}
        \captionsetup{font={small,bf,stretch=1}}
        \label{table1}
    \end{table}


    例2：带斜线的表
    \begin{table}[H]
        \centering
        \setlength{\tabcolsep}{4mm}
        \begin{threeparttable}  
            \begin{tabular}{c|ccccc}  
                \hline
                \diagbox{数值格式}{所需时间(s)}{时刻} & $t=20$ &$t=40$ &$t=60$&$t=70$&$t=80$\\ \hline
                PAVF格式 & $6.4571$ &$17.8073$ &$31.6309$&$38.1617$&$43.9885$\\ \hline
                AVF格式 & $19.0332$ &$48.5705$ &$78.6096$&$91.5572$&$122.6046$\\ \hline
            \end{tabular} 
        \end{threeparttable}
        \caption{\centering{带斜线的表}}
        \captionsetup{font={small,bf,stretch=1}}
        \label{table2}%label一定要写在caption的后面.
    \end{table}

    例3：带注释的表
    \begin{table}[H]
        \centering
        \setlength{\tabcolsep}{4mm}
        \begin{threeparttable}  
            \begin{tabular}{c|ccccc}  
                \hline
                \diagbox{数值格式\tnote{*}}{所需时间(s)}{时刻} & $t=20$ &$t=40$ &$t=60$&$t=70$&$t=80$\\ \hline
                PAVF格式 & $6.4571$ &$17.8073$ &$31.6309$&$38.1617$&$43.9885$\\ \hline
                AVF格式 & $19.0332$ &$48.5705$ &$78.6096$&$91.5572$&$122.6046$\\ \hline
            \end{tabular} 
            \captionsetup{font={small,bf,stretch=1}}
            \begin{tablenotes}    %这行要添加， 从这开始
                \footnotesize               %这行要添加
                \item[*]  统一起见，两种格式均采用迭代算法。     %这行要添加     。
            \end{tablenotes}            %这行要添加
        \end{threeparttable}
        \caption{\centering{带注释的表}}
        \label{table3}
    \end{table}
    更复杂的表格参见网络版本。

    引用表格依然可以采用ref命令。例如：表$\ref{table1}$很简单，表$\ref{table2}$困难一点，表$\ref{table3}$更困难一点。









    \newpage
    \section{算法实现}
    本节写一下算法，算是为计算机和数学的同学多写一点。至于其他的化学流程图，物理动态图请参见网络。
    %此命令放在正文，紧跟\section{}命令，手动为页眉添加节标题
    \markboth{\text{第}\thesection\text{章} \quad 算法实现}{\text{第}\thesection\text{章} \quad 算法实现}  

    \begin{algorithm}[!ht]
        \renewcommand{\algorithmicrequire}{\textbf{输入：}}
        \renewcommand{\algorithmicensure}{\textbf{输出：}}
        \caption{算法的名称}
        \label{power0}
        \begin{algorithmic}[1] % 控制是否有序号
            \REQUIRE  给定初值% input 的内容
            \ENSURE 结果% output 的内容
            
            \STATE 系数矩阵;
            \STATE 右端项; 
            
             % for loop
            \FORALL {$k=1,2,\cdots,Iter$}
                \STATE $N=N(\widehat{\varPhi}^0,\widehat{\varPhi}^1)$;
                \STATE $A=A(\widehat{\varPhi}^1)$;
                \STATE $\widehat{\varPhi}^2=A\backslash N$;
                \STATE $\widehat{\varPhi}^0=\widehat{\varPhi}^1$;
                \STATE $\widehat{\varPhi}^1=\widehat{\varPhi}^2$;
            \ENDFOR
            
             %if ... else
            \IF {$a < 0$}
                \STATE $y = 2$;
            \ELSE
                \STATE $y = 4$;
            \ENDIF
            
            % while 
            \WHILE {a < 0}
                \STATE a--;
            \ENDWHILE
            
             do ... while
            \REPEAT
                \STATE a--
            \UNTIL {a < 0}
            
            \STATE \textbf{return} $\widehat{\varPhi}^1$.
        \end{algorithmic}
    \end{algorithm}

    可以看到的是：里面基础的for,while,if,do……while都有，更多需求可参见网络。
    










    \newpage
    %\vspace*{0pt}
    \section{补充说明}
    %此命令放在正文，紧跟\section{}命令，手动为页眉添加节标题
    \markboth{\text{第}\thesection\text{章} \quad 补充说明}{\text{第}\thesection\text{章} \quad 补充说明}  

    下面说一点整篇文章要注意的地方！！！

    \begin{enumerate}[1)]
    \item 在文篇的开篇，也就是文章的封面，是需要使用者使用word将自己的题目，姓名等填写完成后转成pdf，\textcolor{red}{将其重命名为cover.pdf后上传到与main.tex同一目录后使用的}。封面是可以利用latex打出来的，但使用没有这种方式来的方便。
    \item 附录中写代码的Untitled\_dedai.m也需要放在与main.tex同一目录。并且，mcode库是处理matlab代码的，对于python,c,java等语言要另做处理，可参考网络。
    \item 参考文献中的href命令是给文献加超链接的，可以在pdf中直达文献所在位置，方便别人阅读。但如果不需要可以删除。另外，参考文献也可使用bib.tex,那是一种更加规范的格式，但更复杂，可参考网络。
    \item 在更改标题，也就是更改section的时候，后面的markboth中和section相同的也要一并更改，不然页眉无法正常显示。
    \end{enumerate}






    \newpage
    \section*{附录}
    \phantomsection
    \addcontentsline{toc}{section}{附录}
    \fancyhead[C]{\zihao{5}\songti{附录}}
    所推导出的迭代算法代码如下：
    \lstset{
        caption={迭代算法},
        aboveskip=0pt,%解决代码和文本之间的间距问题
    }
    \lstinputlisting[language=Matlab]{Untitled_dedai.m}

    \textcolor{red}{「如果没有代码可以放点别的，实在不有需要可删除。」}



    \newpage
    \section*{参考文献}
    \phantomsection
    \addcontentsline{toc}{section}{参考文献}
    \fancyhead[C]{\zihao{5}\songti{参考文献}}
    \begingroup  % 去掉thebibliography环境自带的“参考文献”标题
    \renewcommand{\section}[2]{} 

    \begin{thebibliography}{99}
        %\addtolength{\itemsep}{-1pt}  % 缩小行距
        \addtolength{\itemsep}{0em} % 缩小参考文献间的垂直间距
        \setlength{\itemsep}{0pt}
        \bibitem{ref1}\href{https://www.sciencedirect.com/science/article/pii/S0893965921003098}{H.G. Lee, A non-iterative and unconditionally energy stable method the Swift–Hohenberg equation with quadratic–cubic nonlinearity[J]. Applied Mathematics Letters, 2021, 123: 107579.}
        \bibitem{ref2}\href{https://www.sciencedirect.com/science/article/pii/S0893965921000847}{Z.F. Weng, Y.F. Deng, Q.Q. Zhuang, S.Y. Zhai, A fast and efficient numerical algorithm for Swift–Hohenberg equation with a nonlocal nonlinearity[J]. Applied Mathematics Letters, 2021, 118: 107170.}

    \end{thebibliography}
    
    \endgroup




    \newpage
    \section*{致谢}
    \phantomsection
    \addcontentsline{toc}{section}{致谢}
    \fancyhead[C]{\zihao{5}\songti{致谢}}
    时光飞逝，但韶华犹在，在这最该努力奋斗的时光，成为了无知者迷茫的沃土，成为了堕落者无尽的彷徨。
    但天下无不散的筵席，不完美的故事也会有终章。

    2023届数学科学学院学子于2023年05月12日完成，感谢您的使用与反馈。如感觉格式和所要求格式有冲突，可进行更改，本科论文格式已上传，一切以指导老师的要求为准。






















































\end{document}
